% =====================================================================
% Appendix to Section 4 (key: ident). Full proofs and computations.
% Sources: research/tracks/model/notes-final.md §§2, 3, 5, 8 and §12; checks c2, c3, c3b, c4, c6, c7, c9b,
% c13, c15, c16 (research/tracks/model/checks/*.out); referee.md and referee_code/r4, r6;
% research/tracks/universal/notes-final.md Lemma U7; research/tracks/experiments/notes-final.md Prop X13,
% E5, E8 and code/results/e5_gold.md, e8_split_l1.md; research/tracks/pa/notes-final.md Prop 5.1.
% =====================================================================

\section{Proofs and computations for \texorpdfstring{\Cref{sec:ident}}{Section 4}}\label{app:ident}

Proofs follow the model track's final notes; steps they leave out are named. The checks are the seeded scripts of \texttt{research/tracks/model/checks/} (c2--c16), each with its output file; the referee's independent scripts (r1--r7) are in \texttt{referee\_code/}. The model track's referee re-did the proofs of \cref{thm:ident:doob}, \cref{cor:ident:prior,cor:ident:deductive}, \cref{prop:ident:rates}, \cref{prop:ident:sep}, \cref{prop:ident:splits}(a), (b), \cref{thm:ident:kl}, \cref{ex:ident:weaker,lem:ident:zerosum,ex:ident:escape}, \cref{thm:ident:limit}(a), \cref{prop:ident:gold}, \cref{prop:ident:splitlzero}, \cref{prop:ident:spare}(a) and \cref{thm:ident:predded} line by line and found no error beyond the corrections recorded in \cref{rem:ident:exact} and \cref{app:ver}. The results added after the referee report (\cref{thm:ident:limit}(b) with its new proof, \cref{prop:ident:splitlone}, \cref{prop:ident:spare}(d), \cref{ex:ident:lonedq}, \cref{prop:ident:proofs}, \cref{rem:ident:nearmiss}) were checked by the track and again while writing this appendix, not by a referee. Nothing is checked in a proof assistant.

% ---------------------------------------------------------------------
\subsection{Consistency and rates}\label{app:ident:doob}

\paragraph{Proof of \Cref{thm:ident:doob}.}
\emph{Step 1 (joint law; L\'evy).} Let $\Pi$ be the law of $(T,X_1,X_2,\dots)$ on $\Hclass\times\mathcal S^\infty$ with $\Pi(T=t,X\in E)=\pi(t)P_t^\infty(E)$, and $M:=\sum_t\pi(t)P_t^\infty$ its data marginal. Since $\Hclass$ is countable, Bayes' rule gives $\pi_n(A)=\Pi(T\in A\mid X_1,\dots,X_n)$ on the $M$-full event where $M(D_n)>0$. By L\'evy's upward theorem (known), for each fixed $A\subseteq\Hclass$, $\pi_n(A)\to\Pi(T\in A\mid X_\infty)$ $\Pi$-a.s.; the convergence event is a data event, so it holds $M$-a.s.
\emph{Step 2 (the data determine the generator).} For $x\in\mathcal S^\infty$ let $F(x)(s):=\lim_nn^{-1}\#\{i\le n:x_i=s\}$, and $F(x):=\bot$ if some limit fails to exist or the limits do not sum to $1$. By the strong law applied to the countably many $s$, $P_t^\infty(F=P_t)=1$ for each $t$ (on that event the limits are the numbers $P_t(s)$, which sum to $1$). Hence $\Pi(P_T=F(X))=1$.
\emph{Step 3 (the limit lives on $\{T:P_T=F(X)\}$).} Choose a version of $\{\Pi(T=t\mid X_\infty)\}_t$ that sums to $1$ a.s. Each indicator $1\{P_t=F(X)\}$ is $X_\infty$-measurable, so $\Ex[1\{P_T=F(X)\}\mid X_\infty]=\sum_t1\{P_t=F(X)\}\Pi(T=t\mid X_\infty)$ by monotone convergence. The left side is $1$ a.s.\ by Step 2. $F(X)$ ranges $\Pi$-a.s.\ over the countable set $\{P_t:t\in\Hclass\}$; apply Step 1 to the countably many sets $A_\nu:=\{t:P_t=\nu\}$ and intersect the full-measure events. Then, $M$-a.s., $\lim_n\pi_n(A_{F(X)})=\Pi(T\in A_{F(X)}\mid X_\infty)=1$.
\emph{Step 4 (transfer to the truth).} $M\ge\pi(T^*)P_{T^*}^\infty$, so every $M$-null event is $P_{T^*}^\infty$-null and Step 3 holds $P_{T^*}^\infty$-a.s. Under $P_{T^*}^\infty$, $F(X)=P_{T^*}$ a.s.\ by Step 2, so $A_{F(X)}=\Cstar$, and $\pi_n(\Cstar)\to1$.
\emph{Pointwise and total variation.} Write $q_T:=\pi(T)1[T\in\Cstar]/\pi(\Cstar)$. By \cref{cor:ident:prior}, $\pi_n(T)=\pi_n(\Cstar)\,q_T$ for $T\in\Cstar$, so
\[\textstyle\sum_T|\pi_n(T)-q_T|=\sum_{T\in\Cstar}q_T\,(1-\pi_n(\Cstar))+\pi_n(\Hclass\setminus\Cstar)=2\,(1-\pi_n(\Cstar))\to0.\]
\qed

The universal track proves the same statement independently (its Lemma U7, total variation via Scheff\'e's lemma); the exact identity above is ours and gives the same conclusion directly.

\paragraph{Proof of \Cref{cor:ident:prior}.} Every $T\in\Cstar$ has $P_T=P_{T^*}$ as a law on $\mathcal S$, so $P_T(D_n)=P_{T^*}(D_n)$ for every sequence, and the posterior ratio is the prior ratio. \qed

\emph{Computed} (c9b; every likelihood by matching; three seeds, $n\le2000$): $\{z+0=z\}$ and its root split with weights $p_f=(0.5,0.3,0.1,0.1)$ agree to $5.68\cdot10^{-14}$ in log-likelihood; an over-general template, a split missing a root and $T^*$ plus the spare $0=S0$ (weight $0.05$, decaying as $0.95^n$) go to $0$. An earlier check (c9) was tautological for the split (referee m5).

\paragraph{Proof of \Cref{cor:ident:deductive}.} Put $A:=\{T:s\in\Th_d(T)\}$. Then $\pi_n(A)=\pi_n(A\cap\Cstar)+\pi_n(A\setminus\Cstar)$. The second term is at most $\pi_n(\Hclass\setminus\Cstar)\to0$. By \cref{cor:ident:prior} the first equals $\pi_n(\Cstar)\pi(A\cap\Cstar)/\pi(\Cstar)$. Both hold on the single event of \cref{thm:ident:doob}, so simultaneously for all $(s,d)$. For the two cases named after the corollary: under $\Lzero$ with full-support $\Qg$ and positive weights, $\operatorname{supp}P_T=\bigcup\inst(T)$, so every $T\in\Cstar$ has $T^*$'s instance union and hence its $\Th_d$; under $\Lone$ with parameters admissible, $\operatorname{supp}P_T=\Th(T)$, so every $T\in\Cstar$ has $\Th(T)=\Th(T^*)$. \qed

\paragraph{Proof of \Cref{prop:ident:rates}.}
(a) $\ln\frac{\pi_n(T^*)}{\pi_n(T)}=\ln\frac{\pi(T^*)}{\pi(T)}+\sum_{i\le n}Y_i$ with $Y_i:=\ln\frac{P_{T^*}(X_i)}{P_T(X_i)}$ i.i.d. Their negative part is integrable, since $(\ln x)^+\le x$ gives $\Ex Y^-=\Ex(\ln\frac{P_T}{P_{T^*}})^+\le\Ex\frac{P_T(X)}{P_{T^*}(X)}\le1$, and $\Ex Y=\KL(P_{T^*}\|P_T)>0$ by Gibbs' inequality, as $P_T\ne P_{T^*}$. The strong law holds with limit in $(0,\infty]$. If $P_T(X_i)=0$ for some $i$, the ratio is $+\infty$ from then on.
(b) $M(D_n)\ge\pi(\Cstar)P_{T^*}(D_n)$. By the chain rule $\ln\frac{P_{T^*}(D_n)}{M(D_n)}=\sum_{i\le n}\ln\frac{P_{T^*}(X_i)}{M_{i-1}(X_i)}$, and the conditional expectation of the $i$-th term given $D_{i-1}$ is $\KL(P_{T^*}\|M_{i-1})$. Take expectations.
(c) $\frac{\pi_n(A)}{\pi_n(T^*)}=\sum_{T\in A}\frac{\pi(T)}{\pi(T^*)}L_T$ with $L_T:=\prod_{i\le n}\frac{P_T(X_i)}{P_{T^*}(X_i)}$. Use $(\sum a_T)^{1/2}\le\sum a_T^{1/2}$ and, by independence, $\Ex L_T^{1/2}=\rho_T^n$. For the probability bound, $\pi_n(A)\le\pi_n(A)/\pi_n(T^*)$ and Markov's inequality applied to the square root. \qed

For the degeneration: if $T=T^*\cup\{\tau\}$ with weight $\varepsilon$ on $\tau$ and $T^*$'s weights scaled by $1-\varepsilon$, then $P_T\ge(1-\varepsilon)P_{T^*}$, so $\rho_T\ge(1-\varepsilon)^{1/2}\sum_sP_{T^*}(s)=(1-\varepsilon)^{1/2}$.

% ---------------------------------------------------------------------
\subsection{Identification of the generator}\label{app:ident:gen}

\paragraph{Proof of \Cref{prop:ident:sep}.} $\Th(T_1)=\Th(T_2)$: $b$ follows from $a$ and $a\to b$ by MP, and $a\to b$ follows from $b$ by an A1 instance and MP. Under $T_2$: $b$ has root $\neg$; every instance of a logical axiom has root $\to$, $\forall$ or $=$, and $b$ is not an axiom of $T_2$; so no one-node tree concludes $b$, and $\mu_{T_2}(b)\le\Pr[\text{the root is a rule node}]=\alpha_{\mathrm r}$. With $Z_{T_2}\ge1-\alpha_{\mathrm r}$ (\cref{lem:model:lone}(b)), $P_{T_2}(b)\le\alpha_{\mathrm r}/(1-\alpha_{\mathrm r})$. Under $T_1$ the one-node tree citing $b$ has probability $\alpha_{\mathrm{ax}}/2$, and $Z_{T_1}\le1$. The same argument holds under $\Ltwo$. \qed

\emph{Computed} (c6; A1--A3 and MP, formulas of size $\le3$, depth $\le2$; rule probabilities for (ax, lg, MP) $(0.6,0.25,0.15)$ in the proved region, $(0.5,0.3,0.2)$ on its boundary, $(0.3,0.3,0.4)$ and $(0.2,0.4,0.4)$ outside): $P_{T_1}(b)=0.16$--$0.35$, $P_{T_2}(b)=0.007$--$0.04$, $\TV=0.16$--$0.35$; the referee's implementation (r6) confirms both bounds. A route to all parameters, not carried out: $P_T(s)$ is a power series in the rule probabilities with nonnegative coefficients, so the difference is real-analytic on the subcritical region and nonzero near $\alpha_{\mathrm r}=0$.

\paragraph{Proof of \Cref{prop:ident:splits}.}
(a) $\operatorname{supp}P_{T,w}=\bigcup\inst(T)$ by \cref{lem:model:grammar}(c) with positive weights, and $\Th_d$ depends only on the instance union (\cref{def:model:calculus}).
(b) \emph{Templates.} Substituting the body into a pattern occurrence $M(\bar y)$ gives $r(N_1(\bar y,\dots),\dots)$; each $N_j$ has a pattern occurrence, because $\bar y$ are distinct bound variables and a binder production adds one more. Other occurrences of $M$ receive the same body with metavariable-free arguments, which $\DT$ allows. So $\tau_r\in\DT$.
\emph{Instances.} $\inst(\tau)=\bigsqcup_r\inst(\tau_r)$: at a pattern occurrence of $M$ an instance shows the root production of $\theta(M)$, and distinct productions give distinct subtrees there (for a hole production $z_i$, the distinct bound variable $y_i$).
\emph{Probabilities.} For $s\in\inst(\tau_r)$ with $\tau$'s matcher $\theta$, the matcher of $\tau_r$ gives the children's bodies of $\theta(M)$, and the root factorisation gives $\Qg(\theta)=p_r\,\Qg_{\tau_r}(s)$. (Holes and bound indices are drawn alike, so a child body under a binder has the law of a body with one more argument.) Hence $w_\tau\Qg_\tau(s)=\sum_rw_\tau p_r\Qg_{\tau_r}(s)$, and the other components are unchanged.
\emph{Under $\Lone$, $\Ltwo$.} A theory-citation node's conclusion has the same law under $(T,w)$ and $(T',w')$, by the $\Lzero$ statement; every other node has the same law, and a tree's validity and conclusion depend only on its nodes' conclusions. So $\mu_{T'}=\mu_T$ and $Z_{T'}=Z_T$.
(c) A renaming permutes the factors of $\Qg(\theta)$; an argument permutation permutes the holes, which a hole-symmetric $\Qg$ ignores. The converse is (a) for single templates; for first-order templates with only $0$-ary term metavariables, AS Prop~C.1 (under its richness condition) turns equal instance sets into $\tau\equiv\tau'$. That proposition assumes one side has only $0$-ary term metavariables (referee m10), so it says nothing about templates such as A1 $=B\to(C\to B)$. \qed

With \emph{learned} weights the likelihoods differ; that case is \cref{prop:ident:splitlzero,prop:ident:splitlone}. The pre-referee sentence ``the likelihood never tells a schema from its split'' holds only at matched fixed weights (referee m14). Track pa's Prop~2.1 is (b) for $\TInd$ with a shared grammar $\Qg(\varphi)=\Qg(f)\Qg(\cdot\mid f)$: an induction datum with motive root $f$ has probability $w\Qg(f)\Qg(\varphi\mid f)=w\Qg(\varphi)$ under both.

\emph{Computed} (c2, Part A): $z+0=z$ against its split over $0,S,+,\cdot$ with root law $(0.5,0.3,0.1,0.1)$; maximal difference $1.08\cdot10^{-19}$ over all 2849 instances with $|t|\le9$, every instance matched against every template.

\paragraph{Proof of \Cref{rem:ident:splitsL1} (experiments Prop X13).} Under $\Qg$ a closed term has root $f$ with probability $q_f$ and, given the root, children drawn independently from the same closed-term law (every node of a closed term has the same context). So ``cite $\tau$ with probability $w_\tau$ and draw the body of $u$ from $\Qg$'' and ``cite $\tau_f$ with probability $w_\tau q_f$ and draw the bodies of $u_1,\dots,u_r$ from $\Qg$'' give the same law on the cited sentence, jointly with the bodies of $\tau$'s other metavariables. Every later chain step depends only on the current sentence and on the set of MP components, which is the same in $T$ and $T'$ because neither $\tau$ nor any $\tau_f$ is an MP component. \qed

% ---------------------------------------------------------------------
\subsection{Identification in the limit and Gold}\label{app:ident:gold}

\paragraph{Proof of \Cref{thm:ident:limit}(a).} By \cref{thm:ident:doob}, $\pi_n(\Cstar)\to1$, and every $T\in\Cstar$ has $\operatorname{supp}P_T=\operatorname{supp}P_{T^*}$. At most one $L$ can have posterior mass above $\frac12$; once $\pi_n(\Cstar)>\frac12$ it is $\operatorname{supp}P_{T^*}$, which is $\bigcup\inst(T^*)$ under $\Lzero$ (\cref{lem:model:grammar}(c)) and $\Th(T^*)$ under $\Lone$ (\cref{lem:model:lone}(c)). \qed

\paragraph{Proof of \Cref{thm:ident:limit}(b).} The proof of \cref{thm:ident:doob}, with the frequency functional replaced by a support functional.
\emph{A countable parameter that the data determine.} Put $U(T,w):=\bigcup\inst(T)$; it takes countably many values. For a sequence $x$ let $F'(x):=\{x_i:i\ge1\}$. For $w$ in the open simplex, $\operatorname{supp}P_{T,w}=\bigcup\inst(T)$ is a countable set of points of positive probability, each of which occurs a.s., and nothing else occurs; so $P^\infty_{T,w}(F'=U(T,w))=1$. $\Dir_T$ charges only the open simplex. Hence $\Pi(U=F'(X))=1$ for the joint law $\Pi$ of $(T,w,X)$.
\emph{Steps 1 and 3.} The posterior given $D_n$ is $\Pi(\cdot\mid D_n)$ (Bayes' formula with the likelihood as a density with respect to the prior). For each of the countably many sets $A_\nu:=\{(T,w):\bigcup\inst(T)=\nu\}$, L\'evy's theorem gives $\pi_n(A_\nu)\to\Pi(A_\nu\mid X_\infty)$ $M$-a.s. Since $F'(X)$ is $X_\infty$-measurable, $1=\Ex[1\{U=F'(X)\}\mid X_\infty]=\sum_\nu1\{F'(X)=\nu\}\,\Pi(A_\nu\mid X_\infty)$ a.s. So, $M$-a.s., $\pi_n(A_{F'(X)})\to1$.
\emph{Step 4, changed.} Let $E$ be the $M$-null event where this fails. Since $M\ge\pi(T^*)\int P^\infty_{T^*,w}(\cdot)\,\Dir_{T^*}(dw)$, we get $\int P^\infty_{T^*,w}(E)\,\Dir_{T^*}(dw)=0$, so $P^\infty_{T^*,w}(E)=0$ for $\Dir_{T^*}$-a.e.\ $w$, which is Lebesgue-a.e.\ $w$ because the Dirichlet density is positive on the open simplex. For such $w^*$, a.s.\ $F'(X)=\bigcup\inst(T^*)$ and $\pi_n(A_{\bigcup\inst(T^*)})\to1$. \qed

The pre-referee proof said that the steps of \cref{thm:ident:doob} ``go through with the countable sum replaced by an integral''. That sentence is withdrawn: Step 3 needs a countable family of sets, which the support functional provides and the law itself does not (referee M4, who suggested this route).

\paragraph{Sketch of \Cref{thm:ident:limit}(b$'$).} $P_{T,w}$ lies in the Polish space of probability vectors on $\mathcal S$ ($\ell^1$ metric), and the frequency functional $F$ recovers it a.s. For a countable convergence-determining family of bounded continuous functions $g$, L\'evy's theorem gives that the posterior mean of $g(P)$ converges to $g(F(X))$ a.s. So the posterior law of $P_{T,w}$ converges weakly to the point mass at $P^*$, which gives the statement for total-variation balls; transfer to Lebesgue-a.e.\ $w^*$ as in (b). Not written out: the choice of the family and the passage from weak convergence to total-variation balls.

\paragraph{Proof of the counter-statement in \Cref{rem:ident:exact}.} Fix $T$. The set $W_T:=\{w:P_{T,w}=P^*\}$ is the intersection of the simplex with an affine subspace, since $w\mapsto P_{T,w}$ is affine. If it had positive Lebesgue measure in the simplex, the affine map would be constant on an open set, hence on the whole simplex, and every component of $T$ would have law $P^*$ (for $|T|=1$ the condition is $\Qg_\tau=P^*$). The posterior on $w$ given $T$ is absolutely continuous with respect to $\Dir_T$. So the posterior mass of the exact class is $0$ at every $n$ unless some template $\sigma$ has $\Qg_\sigma=P^*=\sum_\tau w^*_\tau\Qg_\tau$. For each $\sigma$, linear independence of $\{\Qg_\tau:\tau\in T^*\}$ allows at most one such $w^*$, and there are countably many $\sigma$. \qed

\paragraph{Proof of \Cref{prop:ident:gold}.} (a) is \cref{thm:ident:doob}. (b) Enumerate $L_\infty=\{e_1,e_2,\dots\}$ and build the text in stages. At stage $i$, with $\sigma$ the prefix built so far, pick $k_i>k_{i-1}$ with $\sigma\cup\{e_i\}\subseteq L_{k_i}$ (possible since the chain increases to $L_\infty$) and append $e_i$; the new prefix $\sigma'$ has positive $P_{k_i}$-probability. By (a), applied to $P_{k_i}^\infty$ conditioned on the positive-probability prefix $\sigma'$, almost every continuation reaches $g_n=L_{k_i}$ at some finite $n$; pick such a continuation, with elements in $L_{k_i}\subseteq L_\infty$, and cut it at a time where $g_n=L_{k_i}$. The resulting sequence lists every $e_i$ and only elements of $L_\infty$, and $g_n=L_{k_i}\ne L_\infty$ at infinitely many times. (c) By (a) with $i=\infty$, $P_\infty^\infty$-a.s.\ $g_n=L_\infty$ eventually. \qed

\emph{Computed} (E5; \cref{tab:exp:e5}): on data uniform on $L_5$ the mass of $L_5$ is $2\cdot10^{-6}$ at $n=64$, $0.661$ at $128$ and $1.000$ from $256$ in all 10 seeds ($L_5$ gains $3-\log_25\approx0.68$ bits per datum, repaying 72 prior bits and its Dirichlet cost); on Gold's text against the MAP learner over $\{L_1,\dots,L_{60},L_\infty\}$ the first 14 stages end (740 data), which shows only that these stages end (experiments' referee m10).
\emph{The $\ISigma n$ chain} (model \S5.2, heuristic). Under $\Lone$ with theorem data, $\PA$'s generator charges $\Th(\PA)\setminus\Th(\ISigma n)$, so $\ISigma n$ is eventually rejected on $\PA$-data; and $\PA$'s generator charges $\Th(\PA)\setminus\Th(\ISigma 5)$, so $\PA$ is eventually rejected on $\ISigma5$-data. Both masses presumably are tiny, as they need long derivations with complex induction; if so, the number of data needed is of order $1/P_{\PA}(\Th(\PA)\setminus\Th(\ISigma n))$. The masses were not estimated.

% ---------------------------------------------------------------------
\subsection{Splits}\label{app:ident:mdl}

\paragraph{Proof of \Cref{prop:ident:splitlzero}.} (a) The split's instance sets are disjoint (\cref{prop:ident:splits}(b)), so the labelling sum of \cref{lem:model:dirsum} has one term, and $\Qg_{\tau_f}(s)=\Qg_\tau(s)/p_f$ for $s\in\inst(\tau_f)$. So $\PDir{T'}(D)=\DirMult(n_f;\alpha)\prod_i\Qg_\tau(X_i)/p_{f_i}$, while $P_T(D)=\prod_i\Qg_\tau(X_i)$.
(b) Stirling: $\ln\Gamma(a+m)=(a+m-\frac12)\ln m-m+\frac12\ln2\pi+o(1)$ as $m\to\infty$. Applied to $\ln\Gamma(K\alpha+n)$ and the $K$ terms $\ln\Gamma(\alpha+n_f)$,
\begin{align*}\ln\DirMult(n_f;\alpha)={}&\textstyle\sum_fn_f\ln n_f-n\ln n+(\alpha-\frac12)\sum_f\ln n_f\\&\textstyle{}-(K\alpha-\frac12)\ln n+\frac{K-1}2\ln2\pi+\ln\frac{\Gamma(K\alpha)}{\Gamma(\alpha)^K}+o(1).\end{align*}
With $\ln n_f=\ln n+\ln\hat r_f$ the first two terms are $n\sum_f\hat r_f\ln\hat r_f$, and the $\ln n$ terms collect to $(K\alpha-\frac K2-K\alpha+\frac12)\ln n=-\frac{K-1}2\ln n$. Subtracting $\sum_fn_f\ln p_f$ turns $n\sum_f\hat r_f\ln\hat r_f$ into $n\KL(\hat r\|p)$.
(c) Pearson's and Wilks's $\chi^2$ limit (known), with all $n_f\to\infty$ a.s.\ since $p$ is in the open simplex. (d) The strong law for $\hat r$; $\ln\DirMult(n_f;\alpha)=n\sum_f\hat r_f\ln\hat r_f+O(\ln n)$ also when some $r_f=0$. (e) The prior. \qed

With $\alpha=\frac12$ the term $(\alpha-\frac12)\sum\ln\hat r_f$ vanishes, and in the well-specified case $\Ex\Delta_n\approx\frac{K-1}2-\frac{K-1}2\ln n+c_{K,\alpha}$; for $K=4$, $c_{4,1/2}=\ln\pi^{-2}+\frac32\ln2\pi=0.467$, which gives the predictions of \cref{tab:ident:c2}.

\begin{table}[ht]
\centering\small
\begin{tabular}{@{}lrrrrr@{}}
\toprule
$n$ & $10^2$ & $10^3$ & $10^4$ & $10^5$ & $10^6$\\
\midrule
well specified, mean $\Delta_n$ & $-4.93$ & $-8.34$ & $-11.77$ & $-15.27$ & $-18.85$\\
\quad prediction & $-4.94$ & $-8.39$ & $-11.85$ & $-15.30$ & $-18.76$\\
\quad $P(\Delta_n>0)$ & $0.007$ & $0$ & $0$ & $0$ & $0$\\
misspecified, mean $\Delta_n$ & $-0.06$ & $39.7$ & $482.2$ & $4917.5$ & $49356.3$\\
\quad prediction & $-1.50$ & $39.5$ & $480.4$ & $4920.4$ & $49351.8$\\
\bottomrule
\end{tabular}
\caption{\Cref{prop:ident:splitlzero}, computed (model check c2, Part B): $\Delta_n$ in nats, $K=4$, $\alpha=\frac12$, $\Qg$'s root law $p=(0.5,0.3,0.1,0.1)$, 400 runs per $n$. Well specified: $r=p$. Misspecified: $r=(0.4,0.3,0.2,0.1)$, $\KL(r\|p)=0.0494$ nats; predictions from (b) with $\hat r$ replaced by $r$, plus $\frac{K-1}2$ (the mean of $\chi^2_{K-1}/2$) in the well-specified case; at $n=100$ the misspecified expansion is not yet accurate. The referee (r4) checked (b) for $\alpha=0.3,1,2$ ($K=3$): the residual tends to $0$ like $1/n$.}
\label{tab:ident:c2}
\end{table}

\paragraph{Proof of \Cref{prop:ident:splitlone}.}
(a) Under $(T',w')$ a theory-citation node chooses $\tau_r$ with probability $w'_r$ and draws the children's bodies from $\Qg$. By the root factorisation and $\inst(\tau)=\bigsqcup_r\inst(\tau_r)$, its conclusion has the law of a citation of $\tau$ whose $M$ has root law $w'$. All other nodes have the same laws, and validity and conclusions depend only on the nodes' conclusions; so $\mu$ and $Z$ coincide, under $\Lzero$, $\Lone$ and $\Ltwo$.
(b) By Tonelli,
\[\textstyle\Ex[\mathrm{BF}_{n+1}\mid D_n]=\int\prod_{i\le n}\frac{P^{(p')}}{P^{(p)}}(X_i)\sum_{x:P^{(p)}(x)>0}P^{(p')}(x)\,\Dir(dp')\le\mathrm{BF}_n.\] Ville's inequality (known; IL Lemma~C.4) and the supermartingale convergence theorem (known).
(c) Write
\[\textstyle\ln\mathrm{BF}_n=\ln\frac{\PDir{T'}(D_n)}{P^{(r)}(D_n)}+\sum_{i\le n}\ln\frac{P^{(r)}(X_i)}{P^{(p)}(X_i)}.\]
The second sum divided by $n$ tends to $\KL(P^{(r)}\|P^{(p)})$ a.s., as in \cref{prop:ident:rates}(a). For the first term:
\emph{upper bound:} $\PDir{T'}(D_n)/P^{(r)}(D_n)$ is a nonnegative supermartingale under $P^{(r)}$, as in (b); it converges a.s.\ to a finite limit, so $\limsup_nn^{-1}\ln$ of it is $\le0$.
\emph{Lower bound:} for $\gamma>0$ let $U_\gamma:=\{p':\KL(P^{(r)}\|P^{(p')})<\gamma\}$. By the continuity lemma below it contains a neighbourhood of the interior point $r$, so $\Dir(U_\gamma)>0$; let $\Dir_U$ be the normalised restriction. By Jensen's inequality,
\[\textstyle\ln\int_{U_\gamma}\prod_i\frac{P^{(p')}}{P^{(r)}}(X_i)\,\Dir(dp')\ge\ln\Dir(U_\gamma)+\sum_iY_i,\qquad Y_i:=\int\ln\frac{P^{(p')}(X_i)}{P^{(r)}(X_i)}\,\Dir_U(dp').\]
The $Y_i$ are i.i.d., and by Tonelli on the positive and negative parts
\[\textstyle\Ex Y^+\le\int\Ex\frac{P^{(p')}(X)}{P^{(r)}(X)}\,\Dir_U(dp')\le1,\qquad\Ex Y=-\int\KL(P^{(r)}\|P^{(p')})\,\Dir_U(dp')\ge-\gamma.\] By the strong law, $\liminf_nn^{-1}\ln[\PDir{T'}(D_n)/P^{(r)}(D_n)]\ge-\gamma$ a.s.; let $\gamma\downarrow0$ along a sequence.
\emph{Continuity lemma: $\KL(P^{(r)}\|P^{(p')})\to0$ as $p'\to r$.} For a tree $\pi$, $\Pr_{p'}(\pi)/\Pr_r(\pi)=\prod p'_f/r_f$ over $\pi$'s theory-citation leaves, $f$ the root production chosen there. Let $L(\pi)$ be the number of these leaves, $\varrho:=\min_fp'_f/r_f\le1$ and $\varrho':=\max_fp'_f/r_f\ge1$. Then $\mu_{p'}(x)\ge\sum_{\pi\to x}\varrho^{L(\pi)}\Pr_r(\pi)\ge\mu_r(x)\,\varrho^{\Ex_r[L\mid x]}$ (Jensen, $\Ex_r[\cdot\mid x]$ over the valid trees concluding $x$), and $Z_{p'}\le\Ex_r[\varrho'^L;\text{valid}]$. So
\[\KL(P^{(r)}\|P^{(p')})\le\ln\tfrac1\varrho\cdot\Ex_r[L\mid\text{valid}]+\ln\frac{\Ex_r[\varrho'^L;\text{valid}]}{Z_r}.\]
Here $\Ex_r[L\mid\text{valid}]\le\Ex_r[N]/Z_r<\infty$, with $N$ the number of derivation nodes ($m<1$). And $\Ex_r[\varrho'^L;\text{valid}]\to Z_r$ as $\varrho'\downarrow1$ by dominated convergence, since $L\le N$ and $\Ex_r[e^{\zeta N}]<\infty$ for small $\zeta>0$ (\cref{lem:model:grammar}(d) for the derivation tree). Both terms tend to $0$ as $p'\to r$. Under $\Ltwo$, $L\le2^d$ and the lemma is immediate. \qed

The pre-referee notes marked ``a derivation-based likelihood does not change the MDL finding'' as proved on the strength of the matched-weight tie alone; that was withdrawn in favour of this proposition (referee M3), and (d) is a conjecture.

\begin{table}[ht]
\centering\small
\begin{tabular}{@{}lp{11.5cm}@{}}
\toprule
part & result\\
\midrule
(a) & split with weights $(u,1-u)$ against the schema with root law $u$, $u\in\{0.1,0.37,0.8\}$: maximal difference $1.05\cdot10^{-14}$\\
latency & 88 of 5885 conclusions depend on $u$ through derived trees\\
(b) & $\delta'=0.5,0.2,0.05$: $P(\max_{n\le3000}\mathrm{BF}_n\ge1/\delta')=0.427,0.133,0.036$ (1000 runs, every $n$ checked)\\
(c) & data root law $0.3$, schema $0.6$: $\KL=0.11487$ nats; mean $n^{-1}\ln\mathrm{BF}_n=0.11301$, $0.11460$, $0.11464$ at $n=10^3,10^4,10^5$\\
(d) & mean $\ln\mathrm{BF}_n=-1.765$, $-6.405$ at $n=10^2$, $10^6$ (400 runs); slope $-0.507$ ($-0.498$ on $n\ge10^4$) against $-0.5$\\
\bottomrule
\end{tabular}
\caption{\Cref{prop:ident:splitlone}, computed (model check c13): a propositional $\Ltwo$ model, depth 2, atom $a$, logical axioms A1--A3 instantiated from a grammar truncated to size $\le5$, MP; $\tau=F\to F$ with $F$ rooted at $a$ or $\neg$, $K=2$, root law $p_a=0.6$; 5885 possible conclusions. Citations are latent only weakly here.}
\label{tab:ident:c13}
\end{table}

\paragraph{E8 windows.} In \cref{rem:ident:e8} the two windows ($n=256\ldots4096$: $1.43$, $1.51$ bits per doubling; $1024\ldots4096$: $1.60$, $1.76$) are both correct for their windows; the prior, not included, adds 76.4 bits for whole. \Cref{prop:ident:splitlone}(b) is a time-uniform bound for a one-component schema with fixed weights, a different setting, so E8 is not a test of it.

% ---------------------------------------------------------------------
\subsection{Spare templates}\label{app:ident:spare}

\paragraph{Proof of \Cref{prop:ident:spare}(a).} The labellings of the data by templates of $T^*\cup\{\sigma\}$ are those for $T^*$, with $n_\sigma=0$. For each labelling, the Dirichlet moment of \cref{lem:model:dirsum} for $T^*\cup\{\sigma\}$ divided by that for $T^*$ is $\frac{\Gamma(A+\alpha_\sigma)}{\Gamma(A+\alpha_\sigma+n)}\frac{\Gamma(\alpha_\sigma)}{\Gamma(\alpha_\sigma)}\Big/\frac{\Gamma(A)}{\Gamma(A+n)}$, the same for every labelling; factor it out. Asymptotics: $\Gamma(a+n)/\Gamma(b+n)\sim n^{a-b}$. \qed

For a symmetric prior ($\alpha_\tau=\alpha_\sigma=\alpha$, $A=K\alpha$) this is experiments' Prop~X7(a) and, at $\alpha=\frac12$, pa's Prop~5.1: $\log_2$ of the inverse is $\frac12\log_2n+\log_2\frac{\Gamma(K/2)}{\Gamma((K+1)/2)}+o(1)$ bits. In pa's code, with $K=8$ and the spare $F\to F$ (15 prior bits), this predicts $19.03$ and $19.82$ bits at $n=1000$ and $3000$; pa computed $19.0$ and $19.8$.

\paragraph{Proof of \Cref{prop:ident:spare}(d1).} $P_{T^*\cup\{\sigma\},w}=\sum_{\tau\ne\tau_1}w_\tau\Qg_\tau+(w_{\tau_1}+w_\sigma)\Qg_{\tau_1}$ for every $w$. By the aggregation property of Dirichlet laws, $(w_{\tau_1}+w_\sigma,(w_\tau)_{\tau\ne\tau_1})\sim\Dir(\alpha_1+\alpha_\sigma,(\alpha_\tau)_{\tau\ne\tau_1})$. So $\PDir{T^*\cup\{\sigma\}}$ is the Dirichlet marginal of $T^*$ with $\alpha_1$ replaced by $\alpha_1+\alpha_\sigma$, and with disjoint instance sets the one-labelling formula of \cref{lem:model:dirsum} gives, for every $D_n$ with $n_1$ data in $\inst(\tau_1)$, exactly
\[R_n=\frac{\Gamma(A+\alpha_\sigma)\Gamma(A+n)}{\Gamma(A)\Gamma(A+\alpha_\sigma+n)}\cdot\frac{\Gamma(\alpha_1+\alpha_\sigma+n_1)\Gamma(\alpha_1)}{\Gamma(\alpha_1+\alpha_\sigma)\Gamma(\alpha_1+n_1)}.\]
The limit follows from $\Gamma(a+n)/\Gamma(b+n)\sim n^{a-b}$ and $n_1/n\to w^*_1$ a.s. \qed

We checked the limits in \cref{tab:ident:spare} against this formula: with $A=1$, $\alpha_1=\frac12$, $w^*_1=\frac12$ it gives $\ln R_\infty=\ln(\frac\pi2\cdot2^{-1/2})=0.1050$ at $\alpha_\sigma=\frac12$, $0$ at $\alpha_\sigma=1$ and $\ln\frac23=-0.4055$ at $\alpha_\sigma=2$, as c15 reports.

\paragraph{Sketch of \Cref{prop:ident:spare}(b), (c).} Write $w=((1-u)\tilde w,u)$ with $u\sim\mathrm{Beta}(\alpha_\sigma,A)$ independent of $\tilde w\sim\Dir(\alpha_\tau)$ (aggregation). The per-datum factor is $1-u+u\,\Qg_\sigma(X)/P_{T^*,\tilde w}(X)$. In (b), at $\tilde w=w^*$ its mean is $1-uc$, so the likelihood decays like $e^{-ncu}$ in $u$, and $\int u^{\alpha_\sigma-1}e^{-ncu}du\sim\Gamma(\alpha_\sigma)(cn)^{-\alpha_\sigma}$. In (c) the mean is $1$ (zero drift) and the log-likelihood in $u$ is locally $-nIu^2/2+\sqrt n\,Zu$; the integral against $u^{\alpha_\sigma-1}$ scales as $n^{-\alpha_\sigma/2}$. Not written out: uniformity in $\tilde w$ near $w^*$, and almost-sure or in-probability control of the fluctuations. This is a boundary Laplace approximation of the kind analysed for overfitted mixtures by \citet{rousseau2011asymptotic}, whose setting has continuous component parameters (checked by the model track's referee); their result is not applied here.

\paragraph{Sketch of \Cref{prop:ident:spare}(d2).} $T^*$ and $T^*\cup\{\sigma\}$ parametrise the same family of laws, so the two marginals integrate the same likelihood against two prior densities on that family, and Laplace's method gives the ratio of the densities at $P^*$. Not written out: the Laplace step.

\begin{table}[ht]
\centering\small
\begin{tabular}{@{}llp{10.4cm}@{}}
\toprule
case & source & result\\
\midrule
(a) & E5(a) & $0=S0$ never cited, $u\cdot0=0$ never matched: $\log_2$ BF $-1.00$, $-6.83$ at $n=1$, $4096$, equal to the formula on every seed; slope $-0.500$\\
(a), (b) & c3, r4 & formula matched to 4 decimals; slopes $-0.499$, $-0.498$; referee: $-\alpha_\sigma$ to 3 decimals for $\alpha_\sigma=\frac12,1,2$\\
(c) & E5(a2) & nested $\varphi(Sz)$ in $\{\varphi(z)\}$: mean $\log_2$ BF $-2.196$ at $n=10^2$, $-7.200$ at $10^8$ (1000 draws per $n$); slope $-0.251\pm0.002$\\
(c) & c3b, r4 & $-0.265$ on $10^3\ldots10^7$ (s.e.\ $\approx0.07$ per point); referee $-0.230$, $-0.523$, $-0.961$ for $\alpha_\sigma=\frac12,1,2$\\
(d1) & c15 & formula equals the numerical integral to 5 decimals; at $n=10^6$ within $0.003$ of the limit\\
(d2) & c15 & $w^*_1=0.3$: $\ln R_n$ changes by $-0.031$, $-0.021$, $-0.004$ per decade; $w^*_1=\frac12$ ($P^*=\Qg_\sigma$, density infinite): $+0.22$, $+0.20$, $+0.13$\\
\bottomrule
\end{tabular}
\caption{Spare templates, computed. Slopes are of the mean log Bayes factor against $\log n$ and so do not depend on the base. The primary estimate for (c) is E5(a2); its rate remains a proof sketch.}
\label{tab:ident:spare}
\end{table}

% ---------------------------------------------------------------------
\subsection{Misspecification}\label{app:ident:misspec}

\paragraph{Proof of \Cref{thm:ident:kl}.} On $\operatorname{supp}P_H$, $\ln(P_T/P_{T_0})=Y_0-Y_T$ with $Y_T:=\ln(P_H/P_T)$. $Y_T$ has an integrable negative part, $\Ex(\ln\frac{P_T}{P_H})^+\le\Ex\frac{P_T}{P_H}\le1$, and mean $K_T$; $Y_0\in L^1$ because also $K_{T_0}<\infty$. By the strong law, $n^{-1}\sum_{i\le n}(Y_{0,i}-Y_{T,i})\to K_{T_0}-K_T$ a.s., the value $-\infty$ allowed when $K_T=\infty$. $\Hclass$ is finite. \qed

Inside $\mathbb M$, members with equal KL but different laws trade places like a zero-drift random walk.

\paragraph{Proof of \Cref{ex:ident:weaker}.} $P_H=P_\sigma$, so \cref{thm:ident:doob} applies with truth $\sigma$. Every $T\in C_\sigma$ has $\bigcup\inst(T)=\inst(\sigma)=\Ic(0+x=x)$ by \cref{prop:ident:splits}(a) (parameters not admissible), hence $\Th(T)=\Th(\Ic(0+x=x))$. Take the structure $\N\cup\{a\}$, standard on $\N$, with $0+a:=0$ and the other operations extended arbitrarily, for instance $Sa:=a$ and $a\circ y:=y\circ a:=a$ for $\circ\in\{+,\cdot\}$ except $0+a=0$. Closed terms denote standard numbers, so every closed instance holds; $\forall x(0+x=x)$ fails at $a$. By \cref{cor:ident:deductive} the posterior probability that $T\vdash\forall x(0+x=x)$ tends to $0$. \qed

\paragraph{Proof of \Cref{lem:ident:zerosum}.} Let $\tau$ cover some $t=0\in Z$.
\emph{Shape.} The datum has no binders, so $\tau$'s rigid skeleton has none. A pattern occurrence of positive arity needs bound variables in scope, so every metavariable is $0$-ary. The only formula position in $t=0$ is the root.
\emph{Root metavariable.} Then $\tau$ is a formula metavariable, with the false instance $0=S0$.
\emph{Otherwise} $\tau=(u=v)$ with term templates $u,v$. $u$ matches $t$, so its rigid nodes are $+$ and $0$ and its metavariable occurrences are leaves; hence $\mathrm{val}(u\theta)=\sum_jc_j\,\mathrm{val}(\theta(z_j))$, $c_j$ the multiplicity of $z_j$ in $u$. $v$ matches $0$, so $v$ is $0$ or a metavariable $z$.
\emph{$v=0$.} Soundness requires $\sum_jc_j\mathrm{val}(\theta(z_j))=0$ for all $\theta$; $\theta(z_j):=S0$ for all $j$ shows $c_j=0$. So $\tau$ is ground and covers one sentence.
\emph{$v=z$, $z$ not in $u$.} The value of $u\theta$ does not depend on $\theta(z)$; one of $\theta(z)\in\{0,S0\}$ gives a false instance.
\emph{$v=z$, $z$ in $u$.} Soundness requires $\sum_jc_j\mathrm{val}(\theta(z_j))=\mathrm{val}(\theta(z))$ for all $\theta$; unit test values give $c_z=1$ and $c_j=0$ for $j\ne z$. So $u$ contains $z$ once and is otherwise ground; covering $t=0$ forces $\theta(z)=0$, and $\tau$ covers exactly one member of $Z$. \qed

\paragraph{Proof of \Cref{ex:ident:escape}.} By \cref{lem:ident:zerosum} each $T_j$ covers at most $|T_j|$ members of $Z$. Since $\operatorname{supp}P_H$ is infinite, $P_H(Z\setminus\bigcup_j\bigcup\inst(T_j))>0$, so a.s.\ some datum falls outside every $T_j$, which then has likelihood $0$ under $\Lzero$. $P_{\Tesc}>0$ on $Z$ (full-support $\Qg$, $z:=t$). So $\pi_n(\Tesc)=1$ from that datum on. $\Tesc\vdash S0=0$ (instance $z:=S0$); $\Q\vdash\neg S0=0$ (Q1). $\Q$'s axioms are sentences outside $Z$, so a theory containing only them has likelihood $0$ at the first datum. \qed

\emph{Computed} (c4, Part 2): twenty sound theories (ground equations for the 20 most likely data, which carry mass $0.956$, weighted as the human's conditional law) share prior $0.99$ against $\Tesc$; $S0=0$ reached posterior $\ge0.99$ in 2000/2000 runs, median 17 data, 90.2\% by $n=50$. \Cref{prop:sound:misspec} uses this example for the verifier.

\begin{table}[ht]
\centering\small
\begin{tabular}{@{}lrrrrrr@{}}
\toprule
& \multicolumn{3}{c}{$\lambda=1$} & \multicolumn{3}{c}{$\lambda=2$}\\
\cmidrule(lr){2-4}\cmidrule(lr){5-7}
seed & 11 & 12 & 13 & 11 & 12 & 13\\
\midrule
distinct data at $n=10^4$ & 282 & 325 & 330 & 282 & 325 & 330\\
hybrid $-$ memorisation, $n=10^3$ & $\ge1143$ & $\ge1207$ & $\ge1065$ & $\ge3344$ & $\ge3505$ & $\ge3063$\\
hybrid $-$ memorisation, $n=10^4$ & $\ge7370$ & $\ge8613$ & $\ge8625$ & $\ge21703$ & $\ge24719$ & $\ge24596$\\
escape alone $-$ memorisation, $n=10^4$ & $-4234$ & $-2649$ & $-2719$ & $12871$ & $16289$ & $16120$\\
\bottomrule
\end{tabular}
\caption{\Cref{rem:ident:fullclass}, computed (model check c7): $\log_2$ score (prior plus $\Lzero$ marginal likelihood, $\Dir(\frac12)$) relative to memorising all distinct data. Human data $t=0$, $t$ a $\{0,+\}$-tree ($0$ with probability $0.7$); model $\Qg$ root law $(0.5,0.3,0.1,0.1)$ on $(0,S,+,\cdot)$; prior $2^{-\lambda\ell}$ under a simplified code (3 bits per symbol, 2 per metavariable occurrence, Elias $\gamma$ for the number of templates), not $\pi_\lambda$ of \cref{def:model:prior}. Hybrid: ground templates for data seen at least twice plus $z=0$; its score is a lower bound (one labelling). Escape alone: $\{z=0\}$. Positive favours the unsound theory. At $\lambda=1$ the escape alone loses to memorisation by $n=10^4$, but the hybrid does not.}
\label{tab:ident:c7}
\end{table}

\paragraph{Details of \Cref{ex:ident:lonedq}.} We write $T_{\Q0}:=T_\Q\cup\{0+z=z\}$ and $T_{\mathrm{mix}}:=T_\Q\cup\{z=0\}$. Each node of the $\Leq$ grammar cites a template of $T$ (probability $\alpha_{\mathrm{ax}}$, metavariables from $\Qg$), or applies refl ($t\sim\Qg$), sym, trans (valid iff the middle terms agree), congruence for $S$, or congruence for $+$ (a side term $u\sim\Qg$, left or right). A tree is valid iff its steps are valid and every term lies in $U_N$, the closed terms of size $\le N$. $P_T=\mu_T/Z_T$, with $\mu_T$ the least fixpoint of the grammar's equation. These rules are complete for closed equational consequence (Birkhoff; recalled, not checked against the source \citep{birkhoff1935structure}). A Monte Carlo simulation agrees with the fixpoint ($Z=0.54909$ against $0.54965\pm0.00091$; $|z|\le1.31$ on eight conclusion frequencies). Human law: $t=0$, $t$ a closed $\{0,+\}$-term from a grammar choosing $0$ with probability $1-h$ and $+$ with $h$, conditioned on size $\le N$, $h\in\{0.3,0.45\}$, with and without the datum $0=0$.

\begin{table}[ht]
\centering\small
\begin{tabularx}{\textwidth}{@{}l>{\raggedright\arraybackslash}p{3.9cm}>{\raggedright\arraybackslash}p{2.6cm}>{\raggedright\arraybackslash}X>{\raggedright\arraybackslash}p{1.6cm}@{}}
\toprule
grammar & rule probabilities (ax, refl, sym, trans, cong$_S$, cong$_+$) & $\Qg$ root law $(0,S,+)$ & minimiser & settings\\
\midrule
1 & $(.35,.15,.1,.15,.1,.15)$, $m=0.65$ & $(.55,.2,.25)$ & $\Tesc$ with $0=0$; $T_{\mathrm{mix}}$ without & 8/8 unsound\\
2 & $(.5,.1,.05,.2,.05,.1)$, $m=0.60$ & $(.55,.2,.25)$ & as grammar 1 & 8/8 unsound\\
3 & as grammar 1 & $(.8995,.1,.0005)$ & $T_{\Q0}$ in 5; $T_{\mathrm{mix}}$ in 3 (all with $0=0$) & 5/8 sound\\
\bottomrule
\end{tabularx}

\vspace{4pt}
\begin{tabular}{@{}lrrrrr@{}}
\toprule
mean $\ln P_T(t=0)$ by size of $t$ & 1 & 3 & 5 & 7 & 9\\
\midrule
grammar 1: $T_\Q$ & $-1.88$ & $-1.82$ & $-8.45$ & $-15.52$ & $-22.89$\\
grammar 1: $\Tesc$ & $-0.80$ & $-3.23$ & $-5.22$ & $-7.20$ & $-9.19$\\
grammar 3: $T_\Q$ & $-1.62$ & $-1.57$ & $-7.91$ & $-15.41$ & $-23.33$\\
grammar 3: $T_{\Q0}$ & $-1.63$ & $-1.28$ & $-7.34$ & $-14.55$ & $-22.19$\\
grammar 3: $\Tesc$ & $-0.43$ & $-8.58$ & $-16.30$ & $-24.01$ & $-31.72$\\
\bottomrule
\end{tabular}
\caption{\Cref{ex:ident:lonedq}, computed (model check c16). Top: the cross-entropy minimiser over $\{T_\Q,T_{\Q0},\Tesc,T_{\mathrm{mix}}\}$ in the eight settings of each grammar ($N\in\{7,9\}$, $h\in\{0.3,0.45\}$, with and without $0=0$). Bottom: per-datum log-likelihoods at $N=9$. In grammar 1, $T_\Q$ pays about 7 nats per extra $+$ and $\Tesc$ about 2; e.g.\ at $N=9$, $h=0.45$, data $t\ne0$, the cross-entropies are $7.758$ ($T_\Q$), $7.124$ ($T_{\Q0}$), $4.951$ ($\Tesc$), $4.904$ ($T_{\mathrm{mix}}$) nats, so the odds of $T_\Q$ against $T_{\mathrm{mix}}$ fall by about 2.85 nats per datum. In grammar 3, $\Tesc$ pays $\ln\frac1{p_+p_0}\approx7.7$ nats per $+$ through $\Qg$; $T_\Q$'s increments (6.3, 7.5, 7.9 nats between consecutive sizes) are still rising at size 9, so the large-term limit in grammar 3 is not settled.}
\label{tab:ident:c16}
\end{table}

% ---------------------------------------------------------------------
\subsection{Predictive versus deductive}\label{app:ident:pred}

\paragraph{Proof of \Cref{thm:ident:predded}.} (a) is \cref{prop:ident:rates}(b) with the index shifted; (b) is \cref{cor:ident:deductive} with $d=\infty$, and the two cases are those named after it. (c) $\Ex_{P_{T_1}}\ln\frac{P_{T_1}(X)}{P_{T_2}(X)}=\KL(P_{T_1}\|P_{T_2})$ per datum, summed over the data; Pinsker's inequality $\TV\le(\KL/2)^{1/2}$ (known). For the misspecified remark: $M(D_n)\ge\pi(T)P_T(D_n)$ for every $T$ and every sequence, so the cumulative log loss of $M$ exceeds that of $T$ by at most $\ln(1/\pi(T))$. \qed

\paragraph{Proof of \Cref{prop:ident:proofs}.} (a) A tree's probability is the product of its nodes' probabilities. A theory-citation leaf whose template is not recorded contributes $\alpha_{\mathrm{ax}}\sum_\tau w_\tau\Qg_\tau(\mathrm{concl}(v))=\alpha_{\mathrm{ax}}P^{\Lzero}_T(\mathrm{concl}(v))$, the sum over the unrecorded choices factorising over leaves (matching is unique, so $\Qg_\tau(\mathrm{concl}(v))$ is the sum over substitutions). No other node depends on $T$. Validity depends only on conclusions, so $Z_T=\Pr[\text{valid}]$ is a function of the law of theory-leaf conclusions, i.e.\ of $P^{\Lzero}_T$. (b) ``If'' is (a). ``Only if'': the one-node tree citing $s$ as a theory axiom has probability $\alpha_{\mathrm{ax}}P^{\Lzero}_T(s)/Z_T$; equality for all $s$ makes $P^{\Lzero}_T$ and $P^{\Lzero}_{T'}$ proportional, and both sum to $1$. The rest is \cref{thm:ident:doob}, \cref{cor:ident:prior} and \cref{prop:ident:splits}(a), (b). (c) \cref{prop:ident:sep}. \qed

If the template name is recorded at each citation leaf, the argument is easier.

\paragraph{Proof of \Cref{rem:ident:nearmiss}.} $P^K_T$ is a probability on $\mathcal S$ whose normaliser does not depend on the data. The proofs of \cref{thm:ident:doob}, \cref{cor:ident:deductive}, \cref{thm:sound:fixed} and \cref{lem:model:size} use only that each hypothesis carries a fixed law on a countable set and that the data are i.i.d.\ from one of them; they therefore hold with $P^K_T$ in place of $P_T$ when the data are i.i.d.\ from $P^K_{T^*}$. \qed
